\documentclass[10pt]{article}

\usepackage[utf8]{inputenc}

\usepackage[T1]{fontenc}

\usepackage{amsmath}

\usepackage{amsfonts}

\usepackage{amssymb}

\usepackage[version=4]{mhchem}

\usepackage{stmaryrd}

\usepackage{bbold}

\usepackage{graphicx}

\usepackage[export]{adjustbox}

\graphicspath{ {./images/} }



\begin{document}

существование и регулярность почти всюду (топологическая локально евклидовость и аналитичность) компакта $X_{0}$, реализующего минимум меры $\Lambda^{k}$ по всем компактам $X$ с заданной границей $L$. Впоследствии эти теоремы были распространены на случай поверхностей $X$ в римановом пространстве.



Другие предложенные обобщения П. з., в частности обобщение в терминах интегральных потоков, являются в нек-ром смысле эквивалентными постановке П. з. в терминах гомологий.



В классич. Постановке Плато многолерная задача решена в 1969; а именно, доказана теорема: если в римановом пространстве $V^{n}$ задано $(k-1)$-мерное подмногообразие $\Gamma, k \geqslant 3$, то существует поверхность, реализующая минимум хаусдорфовой меры $\Lambda^{k}$ среди всех параметризуемых поверхностей $X$, являюпихся в $V^{n}$ непрерывными $f$-образами $k$-мерных гладких многообразий $M$ с краем, гомеоморфным $\Gamma$ при отображении $f: M \rightarrow V^{n}$.



Наряду с вопросами разрешимости П. з. интересными являются также вопросы единственности и регулярности ее решения. Наиболее продвинуто исследование вопроса о регулярности решения. Показано, что решение П. з., данное Дугласом в $\mathbb{R}^{3}$, не содержит внутренних точек ветвления. Для случая многомерной П. 3. доказана регуля рность решения почти всюду, причем возможность наличия нерегулярных точек подтверждается примерами. Что касается единственности решения П. з., то здесь известны лишь нек-рые достаточные признаки единственности (напр., решениө П. з. единственно, если заданный контур $\Gamma$ имеет однозначную выпуклую проекцию при центральном или параллельном проектировании на нек-рую плоскость). Чтобы подчеркнуть сложность этого вопроса, достаточно сказать, что өсть основания ожидать существование спрямляемых контуров, стягивающих континуум м. п. типа круга.



Jum.: [1] Encyklopädie der mathematischen Wissenschaften, Bd 3, H. 2/3, Lpz., 1903; [2] D a r b o u x G., Leçons sur la théorie générale des surfaces..., 2 éd., pt. $1, \mathrm{P} ., 1914 ;$ [3] B i a nc h i L., Vorlesungen über Differentialgeometrie, 2 Aufl., Lpz.ражения и минимальные поверхности, пер. с англ., М., 1953 [5] M o r r e y C., "Ann. Math.", 1948, v. 49, № 4, p. $807-51$; [6] $\mathbf{R}$ a d o T., On the problem of Plateau, B., 1933; [7] Н и ч e И. С. С., "Математика», 1967, т. $11, N 23$, с. $37-100 ;$; 18$]$ O с с е рм а н P., там же, 1971, т. 15, № 2, с. $104-25 ;[9] \boldsymbol{\Phi}$ о м е H-

к о А. Т., "Изв. АН СССР. Сер. матем.», 1972 , т. 36, с. $1049-104$



\includegraphics[max width=\textwidth, center]{2023_07_08_31aba6443d8c04e652c1g-1}

p. $313-33 ;$ [12] O с с е p м a н P., «Успехи матем. наук», 1967 ,

т. 22, в. 4, c. $55-136$; [13] F e d e r e r H. Geometric measure theory, B. - [u. a.], $1969 ;[14]$ M o r r e y C.. Multiple integral in the calculus of variations, B. - [u. a.], 1966. И. Х. Сабитов.



ПЛАТО МНОГОМЕРНАЯ ЗАДАЧА - термин, обозначающий серию задач, связанных с изучением экстремалей и глобальных минимумов функционала $k$-мерного объема $\operatorname{vol}_{k}$, определенного на $k$-мерных обобщенных поверхностях, вложенных в $n$-мерное риманово пространство $M^{n}$ и удовлетворяющих тем или иным граничным условиям.



В истории развития указанной вариационной задачи (см. Плато задача) выделяются несколько периодов, характеризуюгцхся различными подходами к понятиям «поверхности», «границы», «минимизации» и, соответственно, методами получения минимального решения. М но го ме рная 3 а да ч а Пл а то формулируется так. Пусть $A^{k-1} \subset M^{n}-$ фиксированное замкнутое гладкое $(k-1)$-мерное подмногообразие в римановом пространстве $M^{n}$ и пусть $X(A)$ - класс всех таких пленок (поверхностей) $X \subset M^{n}$, имеющих границей многообразие $A$, причем каждая пленка $X \in X(A)$ допускает непрерывную параметризацию (представима в виде образа нек-рого многообразия с краем), т. е. $X=f(W)$, где $W$ - нек-рое $k$-мерное многообразие с краем $\partial W$, гомеоморфным $A$, а $f: W \rightarrow M^{n}-$ непрерывное отображение, совпадающее с фиксированным гомеоморфизмом на крае $\partial W$, то есть $f: \partial W \cong A$. Вопрос: можно ли найти в классе $X(A)$ пленку $X_{0}$, к-рая была бы в каком-либо разумном смысле минимальной, напр. чтобы ее $k$-мерный объем был наименьшим по сравнению с $k$-мерными объемами других пленок $X$ из этого же класса? Оказалось, что черенөсение классич. "двумерных" методов на многомерный случай наталкивается на серьезные трудности. Это привело к тому, что классич. постановка П. м. з. была на время оставлена и задача была сформулирована в иных (гомологических) терминах. Если отбросить понятие многообразия-пленки с краем $\partial W=A$ и сильно расширить понятие пленки и ее границы, ослабив связь пленки с ее границей (в частности, если рассматривать непараметризованные пленки), отбросив условие $X=f(W)$, то П. м. з. может быть сформулирована на языке обычных целочисленных гомологий $H_{*}$ : найти минимальную пленку $X_{0}$, аннулирующую фундаментальный цикл $[A] \in H_{k-1}(A)$ многообразия $A$ (в предположении, что $A$ ориентируемо), т. е. $i_{*}[A]=0$, $i_{*}: H_{k-1}(A) \rightarrow H_{k-1}(X)$, где $i_{*}$ - гомоморфизм, индуцированный вложением $i: A \rightarrow X$. Для решения П. м. 3. в этой новой расширенной постановке был разработан (см. [1], [2]) геометрич. подход, при к-ром минимизировался функционал $k$-мерной хаусдорфовой меры (объема), определенный на $k$-мерных измеримых компактах (поверхностях) в $M^{n}$, и развита (см. [3], [4]) теория интегральных потоков и варифолдов, носителями к-рых являются $k$-спрямляемые подмножества в $M^{n}$. Однако указанное расширение понятия границы пленки в терминах гомологий $H_{*}$ (то есть $k$-1-мерное многообразие $A$ считается границей $k$ мерной пленки $W$, если при вложении $A \rightarrow W$ фундаментальный цикл $[A]$ аннулируется) означает отход классич. постановки П. м. з., так как, располагая теоремой существования минимального решения в гомологич. классе $X(A)$, по-прежнему ничего нельзя сказать о существовании минимального решения в классе всех пленок, являющихся непрерывными образами многообразий с краем, т. е. допускающих параметризацию. Дело в том, что если многообразие $A$ гомологично нулю (как цикл) в пленке $X_{0}$, то $X_{0}$ не обязательно допускает представление в виде $X_{0}=f\left(W_{0}\right)$, где $W_{0}$ - нек-рое $k$-мерное многообразие с краем.



В классич. постановке, т. е. в терминах пленок вида $X=f(W)$, где $W$ суть $k$-мерные многообразия с краем $A$, П. м. з. была решена (см. [5], [6]). При әтом было замечено, что классическая П. м. З. допускает эквивалентную формулировку на языке бордизмов. Пусть $V$ есть $(k-1)$-мерное компактное ориентированное замкнутое многообразие, $f: V \rightarrow M^{n}-$ непрерывное отображение; пара $(V, f)$ наз. с и нг у л я р ны м б о р д и з мо м $\boldsymbol{M}^{n}$. Два бордизма $\left(V_{1}, f_{1}\right)$ и $\left(V_{2}, f_{2}\right)$ наз. ә к в и в а л е н т ны м и, если существует $k$ мерное ориентированное многообразие $W$ с краем $\partial W=V_{1} \cup\left(-V_{2}\right) \quad$ (где $-V_{2}$ означает $V_{2}$ с обратной ориентацией) и непрерывное отображение $F: W \rightarrow M^{n}$ такое, что $\left.F\right|_{V_{1}}=f_{1},\left.F\right|_{V_{2}}=f_{2}$. Бордизм $(V, f)$ эквивалентен нулю, если $V=\partial W=V_{1}, V_{2}=\phi$. Классы эквивалентности сингулярных бордизмов образуют абелевы группы, к-рые после выполнения процесса стабилизации образуют одну из обобщенных теорий гомологий (теорию бордизмов). П. м. З. формулируется (на этом языке) так: а) можно ли среди всех пленок $X \subset M, A \subset X$ и обладающих тем свойством, что сингулярный бордизм $(A, i)$ (где $i: A \rightarrow X-$ вложение) эквивалентен нулю в $X$, найти $X_{0}$ с наименьпим объемом $\operatorname{vol}_{k} X_{0}$ ? б) Можно ли среди всех сингулярных бордизмов $(V, g)$ (где $g: V \rightarrow M^{n}$ ), әквивалентных данному бордизму $\left(V^{\prime}, g^{\prime}\right)$, найти такой бордизм $\left(V_{0}, g_{0}\right)$, чтобы объем





\end{document}